\documentclass[10pt,a4paper]{amsart}
\usepackage[margin=19mm]{geometry}
\usepackage{amsmath,amssymb,mathtools}
\usepackage[scr=rsfs,cal=euler]{mathalfa}
\usepackage{xfrac}
\usepackage[unicode,hidelinks,bookmarks=false]{hyperref}
\newcommand{\ie}{i.e.\ }
\newcommand{\cf}{cf.\ }
\newcommand{\resp}{respectively\ }
\newcommand{\Subsection}{\subsection}
\providecommand{\nobreakdash}{\nobreak\hbox{-}}
\setlength{\emergencystretch}{2em}
\multlinegap0pt
\usepackage{amssymb}
\usepackage{xcolor}
\usepackage{algorithm}\usepackage{algpseudocode}
\usepackage{braket}
\usepackage{mathtools}
\usepackage[shortlabels]{enumitem}
\usepackage{dsfont}
\newtheorem{theorem}{Theorem}
\usepackage{cleveref}
\crefname{theorem}{Theorem}{Theorems}
\title{Equivalence of mutually unbiased bases via orbits: general theory and a $d=4$ case study}
\author{Amit Te'eni, Eliahu Cohen}
\date{Source: arXiv:2508.15412v1 (CC BY 4.0)}
\begin{document}
\maketitle

\noindent\emph{Source and licence.} Equivalence of mutually unbiased bases via orbits: general theory and a $d=4$ case study. Version arXiv:2508.15412v1 (CC BY 4.0), distributed by arXiv under the Creative Commons Attribution 4.0 International licence, http://creativecommons.org/licenses/by/4.0/, as recorded in the arXiv licence metadata for this exact version and shown on the abstract page https://arxiv.org/abs/2508.15412v1. Full licence notice: \texttt{licenses/arxiv-2508.15412-CC-BY-4.0.txt}.\par\smallskip

\noindent\textit{Editorial scope.} Counted unit: the article's theorem thm:main (source0.tex lines 397--400). The article's complete original proof of it is delivered with it (source0.tex lines 401--411, one \texttt{proof} environment of 11 lines). No mathematics is written, repaired or re-derived: the statement and the proof are the article's own lines, in the article's own notation, and no retained line is substituted or re-flowed.  Retained uncounted as the source-local context the proof needs: lines 304--304.  Nothing was omitted from any retained span, so the package carries no omission record.  No retained line contains a citation, so the wrapper carries no bibliography and no bibliography entry.  No retained line contains a figure environment, an image-inclusion command or drawing code, so no image is delivered and the package declares no asset dependency.  Editorial formatting only, in addition: an AMS \texttt{amsart} preamble that restates the article's own declarations and macros used by the retained lines, namely the environments \texttt{theorem} on separate counters, together with the standard packages those lines need.  The retained environments print in this document as Theorem 1 in order of appearance, whereas the article prints the same items with the numbering of its own full document.  The complete native source of the article as distributed in the arXiv e-print for this exact version is bundled unchanged as \texttt{sources/183625/source0.tex}.
\subsection{Subsets of $q$-unbiased bases}\label{sec:Vq}

\begin{theorem}\label{thm:main}
	If $ \left( q_1, \ldots, q_{k-1} \right) $ is a MUB list, then $ \mathcal{N} \left( q_1, \ldots, q_{k-1} \right) $ is closed under the action of the simultaneous stabilizer  $ \mathrm{U} \left( n \right)_{q_1, \ldots, q_{k-1} } $.
	Moreover, for any $p,r \in \mathcal{N} \left( q_1, \ldots, q_{k-1} \right) $, the two MUB lists $ \left( q_1, \ldots, q_{k-1}, p \right) $ and $ \left( q_1, \ldots, q_{k-1}, r \right) $ are equivalent iff $p$ and $r$ belong to the same $\mathrm{U} \left( n \right)_{q_1, \ldots, q_{k-1} }$-orbit in $ \mathcal{N} \left( q_1, \ldots, q_{k-1} \right) $.
\end{theorem}

\begin{proof}
	First, recall from \Cref{sec:Vq} that $ \mathcal{V} \left( q_i \right) $ is closed under the action of $ \mathrm{U} \left( n \right)_{q_i} $. This implies the intersection $\mathcal{N} \left( q_1, \ldots, q_{k-1} \right) = \bigcap_{i=1}^{k-1} \mathcal{V} \left( q_i \right) $ is preserved by each stabilizer $ \mathrm{U} \left( n \right)_{q_i} $, hence by $\mathrm{U} \left( n \right)_{q_1, \ldots, q_{k-1} }$. Thus, indeed there is a well-defined $\mathrm{U} \left( n \right)_{q_1, \ldots, q_{k-1} }$-action on $ \mathcal{N} \left( q_1, \ldots, q_{k-1} \right) $.

	By definition, the two lists $ \left( q_1, \ldots, q_{k-1}, p \right) $ and $ \left( q_1, \ldots, q_{k-1}, r \right) $ are equivalent iff there exists a unitary $U$ that satisfies both conditions: 
	\begin{itemize}
		\item $ U \cdot q_i = q_i $ for all $1 \leq i \leq k-1$;
		\item $U \cdot p = r$.
	\end{itemize}
	$U$ obeys the first condition iff it stabilizes all $q_i$, which occurs iff $ U \in \mathrm{U} \left( n \right)_{q_1, \ldots, q_{k-1} } $.
	By the definition of orbit, there exists an element $U \in \mathrm{U} \left( n \right)_{q_1, \ldots, q_{k-1} }$ that maps $p$ to $r$ iff these two points belong to the same $\mathrm{U} \left( n \right)_{q_1, \ldots, q_{k-1} }$-orbit. This completes the proof.
\end{proof}

\end{document}
